\documentclass[12pt,letterpaper]{article}

\usepackage[utf8]{inputenc}
\usepackage[T1]{fontenc}
\usepackage[spanish,mexico]{babel}
\usepackage{lmodern}
\usepackage{microtype}
\usepackage[letterpaper,margin=2.5cm,headheight=15pt]{geometry}
\usepackage{xcolor}
\usepackage{graphicx}
\usepackage{booktabs}
\usepackage{tabularx}
\usepackage{array}
\usepackage{float}
\usepackage{fancyhdr}
\usepackage{titlesec}
\usepackage{enumitem}
\usepackage{amsmath}
\usepackage{hyperref}
\usepackage{caption}

\definecolor{navy}{HTML}{17324D}
\definecolor{blue}{HTML}{2878A5}
\definecolor{pale}{HTML}{EAF2F7}
\definecolor{ink}{HTML}{263746}
\definecolor{muted}{HTML}{566A78}

\hypersetup{
  colorlinks=true,
  linkcolor=navy,
  urlcolor=blue,
  citecolor=navy,
  pdftitle={Comparación de estrategias de búsqueda y optimización},
  pdfauthor={Torres Flores Raynier Jetsahe}
}

\setlength{\parindent}{0pt}
\setlength{\parskip}{0.55em}
\setlist[itemize]{leftmargin=1.5em,itemsep=0.25em,topsep=0.35em}
\renewcommand{\arraystretch}{1.2}
\titleformat{\section}{\Large\bfseries\color{navy}}{\thesection.}{0.6em}{}
\titleformat{\subsection}{\large\bfseries\color{blue}}{\thesubsection}{0.6em}{}
\titlespacing*{\section}{0pt}{2.0em}{0.7em}
\titlespacing*{\subsection}{0pt}{1.2em}{0.4em}

\pagestyle{fancy}
\fancyhf{}
\fancyhead[L]{\small Inteligencia Artificial}
\fancyhead[R]{\small Informe experimental}
\fancyfoot[L]{\small Torres Flores Raynier Jetsahe}
\fancyfoot[R]{\thepage}
\renewcommand{\headrulewidth}{0.4pt}
\renewcommand{\footrulewidth}{0.4pt}

\newcommand{\metric}[1]{\textcolor{navy}{\textbf{#1}}}
\newcommand{\resultfigure}[4]{%
  \begin{figure}[H]
    \centering
    \IfFileExists{resultados/#1}{%
      \includegraphics[width=#2\textwidth]{resultados/#1}%
    }{%
      \fcolorbox{blue}{pale}{%
        \parbox[c][4.0cm][c]{0.82\textwidth}{%
          \centering
          Gráfica pendiente de generar.\\
          Ejecute el experimento rápido para crear\\
          \texttt{\detokenize{resultados/#1}}.%
        }%
      }%
    }
    \caption{#3}
    \label{#4}
  \end{figure}
}

\begin{document}

% PORTADA
\begin{titlepage}
  \pagecolor{navy}
  \color{white}
  \centering
  \vspace*{1.0cm}
  {\Large\scshape Informe académico\par}
  \vspace{0.7cm}
  {\color{blue}\rule{0.78\textwidth}{1.2pt}\par}
  \vspace{1.4cm}
  {\Huge\bfseries Comparación de estrategias\\[0.25em]
  de búsqueda y optimización\par}
  \vspace{0.8cm}
  {\Large Informe de experimentación\par}
  \vspace{1.6cm}
  {\large\textbf{Asignatura}\par}
  \vspace{0.2cm}
  {\Large Inteligencia Artificial\par}
  \vspace{1.0cm}
  {\large\textbf{Estudiante}\par}
  \vspace{0.2cm}
  {\Large Torres Flores Raynier Jetsahe\par}
  \vspace{0.8cm}
  {\large\textbf{Grupo}\par}
  \vspace{0.2cm}
  {\Large 1959\par}
  \vfill
  {\color{blue}\rule{0.78\textwidth}{1.2pt}\par}
  \vspace{0.5cm}
  {\large Proyecto: algoritmos de búsqueda y optimización\par}
  \vspace{0.4cm}
  {\large \today\par}
\end{titlepage}
\nopagecolor
\color{ink}
\setcounter{page}{1}

\begin{abstract}
Este informe presenta una comparación experimental de búsqueda sistemática y
metaheurísticas en N-Reinas, mochila 0/1, coloreado de grafos y la función
continua Eggholder. Se explican las estrategias, las representaciones de
solución y la configuración empleada. Los resultados se interpretan separando
dos conceptos: la \emph{factibilidad}, que indica si se respetan las
restricciones, y la \emph{calidad}, que indica qué tan conveniente es una
solución factible. La ejecución documentada usa un presupuesto rápido y una
repetición por instancia; por ello sus resultados son evidencia exploratoria,
no una conclusión estadística ni una garantía de optimalidad.
\end{abstract}

\tableofcontents
\newpage

\section{Introducción}

Resolver un problema de optimización implica explorar un conjunto de
alternativas que puede crecer rápidamente con el tamaño de la instancia. La
búsqueda exacta puede certificar una solución, pero su costo puede ser alto.
Las metaheurísticas exploran soluciones de forma guiada y suelen producir
respuestas útiles con un presupuesto acotado, aunque no necesariamente
certifican el óptimo.

El objetivo de este trabajo es comparar cuatro enfoques implementados en el
proyecto: \emph{backtracking} con heurísticas específicas, recocido simulado
(SA), algoritmo genético (GA) y un híbrido GA+SA. Se aplican a N-Reinas,
mochila 0/1 y coloreado de grafos; SA y GA también se comparan en Eggholder,
una función de optimización continua.

La interpretación no se limita a preguntar qué método obtuvo el número más
favorable. También considera si la respuesta respeta las reglas, cuánta calidad
alcanza, cuánto tarda y si la búsqueda exacta terminó o fue truncada por el
límite de cómputo.

\section{Descripción de las estrategias implementadas}

\subsection{Backtracking con heurísticas}

Backtracking construye una solución parcial y explora decisiones de forma
recursiva. Cuando una decisión ya no puede conducir a una respuesta válida o
mejor que la conocida, se descarta esa rama. La implementación no recorre las
alternativas en orden arbitrario: usa ordenamientos, cotas y propagación
específicos de cada problema. En coloreado busca reducir el número de colores y
puede detenerse al alcanzar el límite de nodos configurado.

Si la búsqueda termina sin agotar el límite, sus reglas de poda permiten
certificar el resultado correspondiente. Si se trunca, el mejor resultado
encontrado sigue siendo una solución candidata, pero no se puede afirmar que
sea óptimo ni que no exista una alternativa mejor.

\subsection{Recocido simulado}

SA parte de una solución y genera vecinos mediante una modificación pequeña.
Acepta siempre una mejora y puede aceptar un empeoramiento con probabilidad
\[
  P(\text{aceptar})=\exp\left(-\frac{\Delta}{T}\right),
\]
donde $\Delta$ es el aumento del costo y $T$ es la temperatura. Al principio,
una temperatura mayor permite explorar; al enfriarse, se favorecen movimientos
que no empeoran la solución. El algoritmo conserva la mejor solución visitada,
aunque los movimientos actuales sean temporales.

\subsection{Algoritmo genético}

GA mantiene una población de soluciones. En cada generación selecciona padres
mediante torneo, recombina sus decisiones, aplica mutación y conserva una
cantidad de élite. La población permite explorar varias regiones; la
recombinación reutiliza componentes prometedores y la mutación introduce
diversidad. Si la población o el número de generaciones son pequeños, el
algoritmo puede converger antes de explorar suficiente espacio.

\subsection{Híbrido GA+SA}

El híbrido usa la mejor solución producida por GA como estado inicial de SA.
La combinación busca equilibrar exploración poblacional y refinamiento local.
Su costo incluye ambas fases. No se espera que siempre supere a ambos métodos:
si GA entrega un punto de partida pobre o SA dispone de pocos pasos, el híbrido
puede conservar esa limitación.

\section{Representación utilizada para cada problema}

\subsection{N-Reinas: permutación}

La solución es una permutación de columnas: el índice representa la fila y el
valor representa la columna ocupada. Esta representación garantiza una reina
por fila y por columna desde el inicio; el objetivo restante es eliminar los
ataques diagonales. SA intercambia dos posiciones, y GA utiliza cruce OX y
mutación por intercambio para conservar la propiedad de permutación.

Esta elección evita gastar evaluaciones en tableros que repiten fila o columna.
Sin embargo, una permutación no garantiza que las diagonales sean válidas. Al
aumentar $N$, el espacio de permutaciones crece y las metaheurísticas con un
presupuesto pequeño pueden terminar con ataques diagonales. En backtracking se
elige la fila con menos columnas disponibles y se priorizan valores que
restringen menos las demás filas.

\subsection{Mochila 0/1: vector binario}

Cada posición del vector corresponde a un objeto: 1 indica que se incluye y 0
que se excluye. Esta representación refleja que los objetos son indivisibles.
SA cambia un bit; GA mezcla decisiones de los padres. Ambas metaheurísticas
reparan las propuestas eliminando artículos de menor razón valor/peso hasta
respetar la capacidad.

La reparación hace que las propuestas sean factibles, pero una eliminación
local puede quitar un objeto que fuera importante en una combinación global.
Por ello una selección válida no necesariamente tiene el mayor valor posible.
Backtracking ordena por razón valor/peso y usa una cota superior fraccional
para evitar explorar ramas que no pueden mejorar el valor incumbente.

\subsection{Coloreado de grafos: vector de colores}

Cada posición representa un vértice y su entero indica el color asignado. Una
asignación es factible únicamente si ningún par de vértices adyacentes comparte
color; los conflictos se cuentan por aristas. El objetivo de calidad es
minimizar el número de colores entre las asignaciones factibles.

La búsqueda comienza desde una coloración voraz DSATUR y ensaya cantidades
menores de colores hasta una cota inferior obtenida mediante una clique. SA
recolorea un vértice por paso; GA combina y muta colores por vértice; el híbrido
aplica GA y refina la asignación con SA. Durante la búsqueda pueden aparecer
asignaciones completas pero infactibles. En grafos grandes, demostrar que no
existe una coloración con un color menos puede ser mucho más costoso que hallar
una coloración válida. Por eso una coloración válida de cuatro colores es un
buen resultado observado, pero no prueba que cuatro sea el mínimo si
backtracking se trunca.

\subsection{Eggholder: coordenadas reales}

La solución es un par $(x,y)$ con ambas coordenadas en $[-512,512]$. SA genera
desplazamientos gaussianos locales cuyo tamaño depende del enfriamiento. GA
utiliza cruza lineal y muta una coordenada con ruido; las propuestas se recortan
al dominio permitido.

Eggholder tiene numerosos valles y picos. Pasos cortos pueden quedar atrapados
en un mínimo local; pasos grandes exploran más, pero pueden saltarse regiones
prometedoras. La población de GA ofrece varias propuestas simultáneas, aunque
no garantiza encontrar el mínimo global. En este problema estar dentro del
dominio significa factibilidad; la calidad depende del valor de la función,
para el cual un número menor es mejor.

\section{Configuración experimental}

Las instancias se comparan bajo semillas reproducibles y los algoritmos se
ejecutan sobre las mismas instancias. La ejecución documentada cubre 20
instancias: tres tamaños de N-Reinas, diez filas de mochila, seis grafos y una
instancia de Eggholder. Cada instancia se ejecutó una vez en modo rápido, lo
que produce 79 registros de algoritmo: cuatro métodos por instancia salvo
Eggholder, donde no se usa backtracking.

\begin{table}[H]
  \centering
  \caption{Instancias utilizadas en la ejecución documentada.}
  \begin{tabularx}{\textwidth}{>{\bfseries}p{3.0cm} X}
    \toprule
    Problema & Instancias \\
    \midrule
    N-Reinas & $N=8$, $N=50$ y $N=100$. \\
    Mochila 0/1 & Filas 0--9 del archivo \texttt{knapsack\_5\_items.csv}; se usan pesos, valores y capacidad. \\
    Coloreado & Grafos Erdős--Rényi reproducidos con NetworkX: $n=50$ y $n=1000$, semillas 101, 202 y 303, con $p=6/(n-1)$. \\
    Eggholder & Dos variables reales en el dominio $[-512,512]^2$. \\
    \bottomrule
  \end{tabularx}
\end{table}

La mochila procede del conjunto público
\href{https://www.kaggle.com/datasets/warcoder/knapsack-problem}
{\emph{warcoder/knapsack-problem}}; el archivo incluido fue verificado mediante
SHA-256 frente al valor consignado en el cuaderno de referencia. Se usan
únicamente pesos, valores y capacidad, no las respuestas incluidas en el CSV.
Los grafos se vuelven a generar con el tamaño, la probabilidad y las semillas
del cuaderno, en lugar de usar aristas de ejemplo.

\section{Parámetros utilizados}

La corrida incluida corresponde al comando rápido, con semilla base 2026:

\begin{table}[H]
  \centering
  \caption{Presupuestos y parámetros de la corrida rápida documentada.}
  \begin{tabularx}{\textwidth}{>{\bfseries}p{3.1cm} p{3.4cm} X}
    \toprule
    Estrategia & Parámetros & Justificación e implicación \\
    \midrule
    SA & 100 iteraciones; $T_0=100$; enfriamiento 0.995; $T_{\min}=10^{-8}$. &
    Presupuesto reducido para ejecutar todas las instancias con rapidez. El enfriamiento gradual permite aceptar algunos empeoramientos al comienzo y estabiliza después; la temperatura inicial común no está calibrada de forma óptima para todas las escalas de costo. \\
    GA & Población 12; 12 generaciones; cruza 0.9; mutación 0.05; elitismo 2; torneo 3. &
    Mantiene una búsqueda poblacional pequeña. Más individuos o generaciones aumentarían la exploración y el costo; el elitismo preserva los mejores individuos y la mutación introduce diversidad. \\
    Backtracking & Límite de 10\,000 nodos. &
    Evita que la búsqueda exacta consuma tiempo excesivo en instancias combinatorias grandes. Alcanzar el tope implica truncamiento y no certifica optimalidad. \\
    Repeticiones & Una por instancia. &
    Permite una validación rápida, pero no mide adecuadamente la variabilidad aleatoria entre ejecuciones. \\
    \bottomrule
  \end{tabularx}
\end{table}

Estos valores son un presupuesto práctico de comparación, no hiperparámetros
óptimos universales. Una repetición sirve para verificar el flujo de ejecución
y observar ejemplos; una evaluación robusta de métodos estocásticos debe usar
múltiples semillas y resumir medias, dispersión y tasas de factibilidad.

\section{Resultados}

\subsection{Distinción entre factibilidad y calidad}

\metric{Factibilidad} responde si la solución respeta las restricciones:
capacidad no excedida en mochila, cero ataques en reinas, cero conflictos en
coloreado y coordenadas dentro del dominio en Eggholder.

\metric{Calidad} compara soluciones factibles según el objetivo del problema:
mayor valor en mochila, menos colores en grafos y menor valor de función en
Eggholder. En N-Reinas, todas las soluciones factibles tienen cero conflictos;
por tanto, una vez factible, la calidad queda empatada. Una solución no
factible puede tener un número de objetivo aparentemente atractivo, pero no
debe contarse como respuesta válida.

La brecha porcentual se compara con la mejor calidad factible observada entre
los métodos de esta corrida. Es una referencia empírica, no un porcentaje de
optimalidad, excepto cuando el algoritmo exacto termina y certifica el óptimo.

\subsection{Resumen observado}

\begin{table}[H]
  \centering
  \caption{Resultados representativos de la ejecución rápida. Los tiempos son de esta máquina y una repetición.}
  \small
  \begin{tabularx}{\textwidth}{>{\bfseries}p{2.5cm} X p{2.7cm} p{2.3cm}}
    \toprule
    Problema & Observación & Factibilidad & Calidad / advertencia \\
    \midrule
    N-Reinas, $N=8$ & Backtracking y SA hallaron solución válida. & GA e híbrido no fueron factibles en esta corrida. & Las soluciones válidas tienen 0 conflictos; la calidad factible empata. \\
    N-Reinas, $N=50,100$ & Backtracking resolvió ambos tamaños; las metaheurísticas no lograron cero conflictos. & Solo backtracking fue factible. & SA/GA tenían 100 pasos o 12 generaciones; el presupuesto fue corto para esta instancia. \\
    Mochila & Backtracking obtuvo el mejor valor en las diez filas; las metaheurísticas empataron en nueve. & Todos los métodos fueron factibles. & En la fila 4 GA obtuvo 43 frente a 44: brecha aproximada de 2.3\%. \\
    Grafos, $n=50$ & Todos los métodos encontraron coloración de cuatro colores. & 100\% factibles en las seis instancias. & El resultado mínimo no queda certificado si la búsqueda exacta se trunca. \\
    Grafos, $n=1000$ & Todos encontraron coloración válida de cuatro colores. & 100\% factibles. & Backtracking alcanzó 10\,000 nodos en las tres semillas; no demuestra que cuatro sea el mínimo. \\
    Eggholder & El híbrido obtuvo $-934.697$; GA $-801.924$; SA $-557.739$. & Los tres permanecieron dentro del dominio. & El híbrido fue el mejor observado; una corrida no permite afirmar superioridad general. \\
    \bottomrule
  \end{tabularx}
\end{table}

Los datos de mochila tienen pocos objetos por instancia, lo que favorece que
backtracking termine rápidamente y encuentre el mejor valor en las diez filas.
En cambio, los tableros grandes de reinas muestran una diferencia marcada
entre hallar un tablero estructuralmente correcto (permutación) y eliminar
todos los ataques diagonales. En Eggholder, el orden observado entre métodos
es compatible con la presencia de mínimos locales y el refinamiento que aporta
SA al punto inicial de GA, pero no constituye prueba estadística.

\section{Visualizaciones}

Las curvas siguientes se generan a partir de la corrida del proyecto y se
incluyen si los PNG existen dentro de la carpeta \texttt{resultados}. Para
compararlas, recuérdese que algunas gráficas representan conflictos, otras
valor y otras colores; la dirección de mejora no es idéntica en todos los
problemas.

\resultfigure{convergence_queens_N_100.png}{0.82}{Convergencia en N-Reinas para $N=100$. Cero conflictos representa una solución factible.}{fig:queens}
\resultfigure{convergence_knapsack_kaggle_opcion1_fila_4.png}{0.82}{Convergencia en la fila 4 de mochila; el valor total se maximiza.}{fig:knapsack}
\resultfigure{convergence_coloring_n_1000_seed_101.png}{0.82}{Convergencia del coloreado en un grafo de 1000 vértices, semilla 101.}{fig:coloring}
\resultfigure{convergence_eggholder_-512_512_2.png}{0.82}{Convergencia en Eggholder dentro del dominio bidimensional.}{fig:eggholder}

\section{Comparación de algoritmos}

\begin{table}[H]
  \centering
  \caption{Comparación cualitativa basada en los mecanismos y resultados observados.}
  \small
  \begin{tabularx}{\textwidth}{>{\bfseries}p{2.8cm} X X}
    \toprule
    Método & Ventaja observada o esperada & Límite importante \\
    \midrule
    Backtracking & Puede certificar soluciones/óptimos al terminar; fue eficaz en mochila y N-Reinas de estos tamaños. & Puede crecer exponencialmente; en grafos de 1000 nodos se truncó en las tres semillas. \\
    Recocido simulado & Movimiento local barato; útil para refinar soluciones y explorar gracias a aceptar temporalmente empeoramientos. & Depende de escala de objetivo, temperatura y presupuesto; no garantiza óptimo global. \\
    Algoritmo genético & Explora varias soluciones a la vez y conserva diversidad mediante mutación. & Población y generaciones pequeñas pueden no bastar; en N-Reinas grandes no encontró cero conflictos en la corrida rápida. \\
    Híbrido GA+SA & Combina recombinación global con refinamiento local; fue el mejor observado en Eggholder. & Acumula el costo de las dos fases y hereda las limitaciones del estado inicial y del presupuesto. \\
    \bottomrule
  \end{tabularx}
\end{table}

Los tiempos absolutos dependen del equipo y no deben generalizarse como una
propiedad intrínseca del algoritmo. La comparación más justa debe mantener
iguales las instancias, semillas y criterios de evaluación, además de reportar
calidad, tasa factible, tiempo y truncamientos juntos.

En los grafos de 1000 nodos, los tiempos de las tres semillas promedian
aproximadamente 17.59\,s para backtracking, 0.41\,s para GA, 0.40\,s para el
híbrido y 0.33\,s para SA. La diferencia se debe interpretar junto con el
objetivo de cada método: las metaheurísticas hallaron cuatro colores válidos
sin certificar que fueran el mínimo; backtracking alcanzó su presupuesto de
10\,000 nodos intentando certificar si era posible usar menos colores. Por
tanto, la ventaja de velocidad observada no implica que las metaheurísticas
hayan demostrado optimalidad.

\section{Discusión}

Los resultados ilustran que el comportamiento depende de la estructura del
problema y de la codificación. La permutación de N-Reinas evita ataques por
fila y columna, pero deja una parte difícil --las diagonales-- para optimizar.
Con solo 100 pasos o 12 generaciones, las metaheurísticas no alcanzaron cero
conflictos en los tableros de 50 y 100 reinas, mientras backtracking sí lo
hizo en esta ejecución.

En mochila, la reparación mantiene todas las propuestas dentro de la
capacidad, lo que explica la factibilidad del 100\%. No obstante, reparar por
razón valor/peso es una decisión local y no garantiza por sí sola el mejor
subconjunto. Como cada instancia incluida contiene pocos objetos, la búsqueda
con poda de backtracking puede resolverlas con rapidez y ofrecer una
referencia más sólida.

En coloreado conviene separar claramente validez y número de colores. Las seis
instancias terminaron con coloraciones válidas de cuatro colores, una señal de
buena calidad observada. Sin embargo, hallar una solución con cuatro colores
solo establece que existe una coloración de ese tamaño; no establece que no
exista otra con tres. El límite de backtracking impidió certificar esa
optimalidad en los grafos grandes.

Eggholder evidencia el efecto del paisaje de optimización. En esta corrida el
híbrido encontró el menor valor de los tres métodos; GA proporcionó una
solución inicial que SA pudo refinar. Aun así, el resultado depende de una
semilla y de presupuestos cortos: se necesitarían varias repeticiones y
presupuestos mayores para evaluar su estabilidad.

Finalmente, la comparación de tiempos debe matizarse. En problemas pequeños,
el costo de inicialización puede ser similar o mayor que el de búsqueda. En
grafos grandes, el trabajo de backtracking crece al intentar certificar una
cantidad menor de colores; una metaheurística puede hallar rápido una
coloración válida, aunque no pruebe que sea óptima. Ningún método es mejor en
todos los criterios y problemas.

\section{Conclusiones}

\begin{enumerate}[leftmargin=1.7em]
  \item La representación influye en qué soluciones puede proponer el
  algoritmo: las permutaciones preservan filas y columnas en N-Reinas, los
  vectores binarios representan decisiones indivisibles en mochila y los
  vectores de color codifican directamente las asignaciones del grafo.
  \item Factibilidad y calidad deben reportarse por separado. Una solución
  factible cumple restricciones, pero puede ser de baja calidad; una solución
  de alta calidad aparente que viola restricciones no es aceptable.
  \item La búsqueda exacta fue adecuada para las instancias pequeñas y para
  N-Reinas hasta $N=100$ en esta corrida, pero alcanzó su presupuesto en los
  grafos grandes. El truncamiento no prueba inviabilidad ni optimalidad.
  \item Las metaheurísticas produjeron resultados útiles en grafos y
  Eggholder, pero los presupuestos rápidos fueron insuficientes para resolver
  todos los tableros de reinas. El híbrido ofreció el mejor valor observado en
  Eggholder, sin que una sola repetición permita generalizar.
  \item Para conclusiones más sólidas se recomienda aumentar las repeticiones,
  conservar semillas, comparar la variabilidad y probar presupuestos
  progresivamente mayores sin ocultar el costo computacional.
\end{enumerate}

\section*{Reproducibilidad}
\addcontentsline{toc}{section}{Reproducibilidad}

Desde la raíz del proyecto, genere nuevamente los resultados y gráficas con:

\begin{verbatim}
.\.venv\Scripts\python.exe -m ia_optimizacion --quick --output resultados
\end{verbatim}

Compile este archivo con pdfLaTeX dos veces para actualizar las referencias y
la tabla de contenidos:

\begin{verbatim}
pdflatex informe_tarea3.tex
pdflatex informe_tarea3.tex
\end{verbatim}

El informe usa las cifras de una ejecución rápida ya guardada en
\texttt{resultados/summary.csv}. Si se regenera la carpeta con otros parámetros
o semillas, actualice también las cifras resumidas de las secciones de
resultados y discusión.

\end{document}
